% !TeX root = ../papers/nc-p3.tex
\section{全四元群両対称型の評価判定と負交換の端点}
\label{sec:negative-exchange-end-values}
\begin{theorem}\label{thm:all-klein-bisymmetric-evaluation}
定理\ref{thm:klein-bisymmetric-irreducible}の全固定色両対称型において、
全次数の標準Koszul左右AS正則性は次の三条件と同値である。
反対階数が6または8であること、辺の六次元関係空間が非零の純テンソルを含まないこと、
左右の全非零部分評価が階数3以上であることである。捻りの場合は逆捻り後に同じ条件を使う。
この判定には薄い係数、角ブロック、点配置の開条件がない。
\end{theorem}
\begin{proof}
定理\ref{thm:bisymmetric-six-or-eight}、系\ref{cor:one-smooth-apolar-as}の必要条件を使う。
反対切断を $\mathrm{Sym}^2 E$ からの写像として読むと、階数6では核が四次元である。
対称四行行列の階数2以下の射影多様体は次元6なので、核の三次元射影空間がこれと交わる。
従って核に二つの一次式の積が入り、補題\ref{lem:bisymmetric-reducible-quick-test}が使える。
階数8では核は二次元である。四文字の平方の像は総文字1により、相手側の $01,23$ の二積しか使わない。
従って核の全体が四つの平方からなる鉛筆である。四列のいずれか二つを同時に消す鉛筆要素があれば
その要素は階数2以下、なければ四列はいずれも非零で相異なる向きを持ち、基底曲線は滑らかである。
後者には系\ref{cor:one-smooth-apolar-as}、前者には同じ補題を用いる。
\end{proof}

\begin{theorem}\label{thm:negative-exchange-minus-one}
$D\ne0,z=-1$ の負交換表示の指定始点の正則商は重み付き射影空間
$[A:B:C:L]\in\mathbb P(1,1,1,2)$ の次の開集合である。
\[
 A+B-C\ne0,\qquad
 (B,(A-C)^2-4L)\ne(0,0),\qquad
 (A,(B-C)^2-4L)\ne(0,0).
\]
商写像は $A=p^2,B=r^2,C=q^2+t^2,L=q^2t^2$ で、一位置移動は $A,B$ を交換する。
従って最小周期は $A=B$ で1、他で2である。
反対階数は一般に12、$A,B$ の片方だけが零なら10、両方零なら8である。
階数8の全点は周期1で、その商は $L/C^2\in\mathbb A^1\setminus\{1/4\}$ である。
この端点と $z\ne0,\pm1$ の正則交換点には相互に同型がない。
\end{theorem}
\begin{proof}
正の交換に戻して八直線による除外を示す。純テンソルは角二ブロックの行列式から
$z\ne1$ だけで除かれる。反対核の対角鉛筆が滑らかまたは可約元を持つことも以前と同様である。
正確には階数が8に達しなくても、平方の四行二列ブロックの行対が比例すれば可約元が生じ、
どの行対も独立なら他の二ブロックの階数2が $z\ne1$ から従い、全階数は8である。
従って評価階数だけを調べればよい。全四座標が非零の候補を消す行列式の符号の矛盾は
$\det K=-\det L=(z-1)/D\ne0$ だけを使う。全低二座標または全高二座標が零なら
厚い評価の階数2に少なくとも一つの薄い式が加わる。低い側または高い側で一座標だけが零なら、
核が二次元のとき他側にも零座標を要することは、定理\ref{thm:exceptional-exchange-eight-lines}と同じ
可逆な二行列の列比較から従う。従って八つの二座標の場合を調べればよい。
座標を $1,s$ とし $sD$ を可逆化して、六行四列の評価の三次小行列式のイデアルを比較すると、
$z=-1$ では八つ全てが同定理の直線と各々の $s^2$ の式になる。
\path{computations/verify_exchange_end_values.py} は等式を有理数上で両方向に検証する。

前二節の頂点復元・軌道証明は、$z=-1$ でも成り立つ。二次式の行列式に使った非零因子は $D,z,z-1$ と
明示した $q,t$、二重二次曲線の除法には $D$ を用いた。二組の座標の全置換を合わせる48語の検査で
追加解を許す因子は $z-1$ だけである。一方の対角差が零の枠の42定数語も全て非零で、
他の八置換が $z'b'=c'$ を強制する条件も同じである。両方零なら $w^2=-1$ で $w\pm1\ne0$、
四角形の28語の比較も使える。このとき $X=-1$ で二つの有理変換の作用が重なるが、
等しい $X,Y$ を持つ表示が $U$ により同型になる逆向きの証明は同じである。
二次曲面では既知の座標は $-2$ となり、この点も先の負成分上の正則閉点として実現する。
いずれの形でも、前節の $z$ または $1/z$、四角形での $X$ または $1/X$、
二次曲面の座標のいずれかが、従来の交換表示との同型を区別する。

$z=-1$ なら $U$ は $r$ だけの符号変更、$US,VS$ は $q,t$ の符号変更、$B$ はこれらを除けば
$q,t$ の交換である。従って射影作用は四座標の任意の符号変更と $q,t$ の交換の全体である。
不変環の射影商が表示の四座標となる。$4D=A+B-C$ であり、八直線のうち四本は
$r=0$, $p\pm q\pm t=0$、他の四本は $p=0$, $r\pm q\pm t=0$ となるから除外式を得る。
$J$ の固定には $A=B$ のほか射影倍率 $-1$ が考えられるが、そのとき $A+B=C=0$ で除かれる。
反対ブロックの色2,3はこの繊維で階数2となり、色0の二小行列はともに $p$、
色1の二小行列はともに $r$ の定数倍である。これが反対階数の式を与える。
\end{proof}

\begin{proposition}\label{prop:zero-exchange-vertical-identification}
$D\ne0,z=0$ の負交換窓は第二長方形の負垂直族と同じ表示である。
必要十分な正則条件は、$(a',b',c',d')=(a,-b,-c,d)$ が
\[
 (a'=b',c'=b'),\quad(d'=b',c'=b'),\quad
 (a'=b',d'=0),\quad(a'=0,d'=b')
\]
の四直線を避けることである。$z=1,D\ne0$ の交換表示には正則点がない。
\end{proposition}
\begin{proof}
$z=0$ でも正の窓の純関係の排除と二座標への帰着に必要な厚い二行列は可逆である。
四行二列の行そのものが零になった場合も、反対核に平方が入る。
八評価イデアルのうち左02,12と右02,03だけに解があり、順に上の四直線と
$s^2=b'$ が得られる。$sD\ne0$ の局所化での計算は同じスクリプトに含めた。
$z=1$ では正表示の $x=(0,0,1,1)$ の部分評価が階数2で、捻りも非正則である。

$z=0$ の負表示の奇位置の文字に1を加え、目標の四位置で順に
$(-,+,+,-),(+,-,+,+),(+,+,-,+),(-,+,+,+)$ の符号を掛けると、
定理\ref{thm:parallel-vertical-orbits}の $t=-1,b_{\rm rect}=c_{\rm rect}=1$、
$(A,B,C,D_{\rm rect})=(a,c,-d,-b)$ のテンソルと全256係数で一致する。
角ブロックの条件は元の $ad-bc\ne0$ に等しい。この恒等式は $B=0$ でも有効である。
従って $B\ne0$ の正則な八代表台の同型と周期には既存の長方形商を適用できる。
$B=0$ なら $a,d\ne0$ により上の四直線のどれにも乗らず、全点が正則である。
\end{proof}
